%Master File:lectures.tex

\lesson{Evolutionary Trees}
\vspace{-2cm}
\begin{center}
 \includegraphics[height=10cm]{additiveGraph}
\end{center}

\keywords{Additive Trees and Parsimony}
%%%%%%%%%%%%%%%%%%%%%%% Next Slide %%%%%%%%%%%%%%%%%%%%%%%
\renewcommand{\Outline}{%
\begin{slide}
\section[1]{Outline}

\begin{minipage}{12cm}
  \begin{enumerate}\squeeze
    \outlineitem{Additive Trees}{addit}
    \outlineitem{Parsimony}{parsimony}
   \end{enumerate}
\end{minipage}\hfill
\begin{minipage}{10cm}
  \vfill
  \includegraphics[width=10cm]{additiveGraph}
  \vfill
\end{minipage}
\end{slide}
\addtocounter{outlineitem}{1}
}

\setcounter{outlineitem}{1}
\Outline
\toptarget{firstoutline}
%%%%%%%%%%%%%%%%%%%%%%% Next Slide %%%%%%%%%%%%%%%%%%%%%%%

\begin{slide}
\section{Recap}
\begin{PauseHighLight}
  \begin{itemize}
  \item We looked at building ultrametric trees which would result if
    mutations accumulated at a regular rate\pause
    \begin{center}
      \includegraphics[width=0.9\linewidth]{ultrametric}
    \end{center}
  \item Unfortunately often our distance are far from ultrametric\pause
  \end{itemize}
\end{PauseHighLight}

\end{slide}
%%%%%%%%%%%%%%%%%%%%%%% Next Slide %%%%%%%%%%%%%%%%%%%%%%%

\begin{slide}
\section[-2]{Additivity of Distances}

\begin{PauseHighLight}
  \begin{itemize}
  \item The condition of ultrametricity is very strong and only works
    with a perfect molecular clock\pause
  \item A less strong condition is \emph{additivity of distances}\pause
    \begin{center}
      \includegraphics[width=0.4\linewidth]{additiveGraph}
    \end{center}
    \begin{align*}
      d_{i\ell} &= d_{ik} + d_{k\ell} & d_{j\ell} &= d_{jk} + d_{k\ell} & 
      d_{ij} &= d_{ik} + d_{jk}\pause\\
      & & d_{k\ell} &= \tfrac{1}{2} (d_{i\ell} + d_{j\ell} - d_{ij})\pause
    \end{align*}
  \end{itemize}
\end{PauseHighLight}

\end{slide}

%%%%%%%%%%%%%%%%%%%%%%% Next Slide %%%%%%%%%%%%%%%%%%%%%%%

\begin{slide}
\section{Additivity and Graphs}

\begin{PauseHighLight}
  \begin{itemize}
  \item Additivity does not assume a constant mutation rate, but only
    that mutations accumulate\pause
  \item That is, the distances are consistent with an evolutionary graph
    \begin{center}
      \includegraphics[width=0.4\linewidth]{additiveGraph}\pause
    \end{center}
  \item This would be true if all mutations were independent of each other\pause
  \end{itemize}
\end{PauseHighLight}

\end{slide}


%%%%%%%%%%%%%%%%%%%%%%% Next Slide %%%%%%%%%%%%%%%%%%%%%%%

\begin{slide}
\section[-2]{Building Additive Trees}

\begin{PauseHighLight}
  \begin{itemize}
  \item For additive distance matrix then the closest nodes are not
    necessarily neighbours\pause
    \begin{center}
      \includegraphics[width=0.5\linewidth]{additiveNeighbour}
    \end{center}
  \item To compute the graph for nodes with an additive distance matrix
    we can compute a new set of ``distances''
    \begin{align*}
      D_{ij} &= d_{ij} - r_i - r_j & r_i &= \frac{1}{|\mathcal{L}|-2}
                                           \sum_{k\in\mathcal{L}} d_{ik}\pause
    \end{align*}
  \end{itemize}
\end{PauseHighLight}

\end{slide}

%%%%%%%%%%%%%%%%%%%%%%% Next Slide %%%%%%%%%%%%%%%%%%%%%%%

\begin{slide}
\section[-2]{Neighbour-Joining Algorithm}

\begin{PauseHighLight}
  \begin{itemize}
  \item From $d_{i,j}$ compute $D_{ij}$\pause
    \begin{center}
      \includegraphics[width=0.4\linewidth]{distanceTable}\hfil
      \includegraphics[width=0.4\linewidth]{D_Table}\hfil
    \end{center}
  \item Choose the pair $(i,j)$ with the smallest $D_{ij}$ value\pause
  \end{itemize}
\end{PauseHighLight}

\end{slide}

%%%%%%%%%%%%%%%%%%%%%%% Next Slide %%%%%%%%%%%%%%%%%%%%%%%

\begin{slide}
\section[-1]{Joining Nodes}

\begin{PauseHighLight}
  \begin{itemize}
  \item Start with a forest of unconnected nodes\pause
  \item We then create a new node $k$ and join $i$ and $j$ to $k$\pause
  \item We compute the distance $d_{km} =\tfrac{1}{2}(d_{im} + d_{jm} -
    d_{ij})$ between $k$ and all other nodes $m$\pause
  \item The distance from $d_{ki} = \tfrac{1}{2}(d_{ij} + r_i - r_j)$
    and $d_{kj} = \tfrac{1}{2}(d_{ij} - r_i + r_j)$\pause
  \item We then remove $i$ and $j$ from our list of nodes to consider
    and add $k$ instead\pause
  \item This is iterated until all nodes have been joined\pause
  \end{itemize}
\end{PauseHighLight}

\end{slide}

%%%%%%%%%%%%%%%%%%%%%%% Next Slide %%%%%%%%%%%%%%%%%%%%%%%

\begin{slide}
\section{Additive Distances}

\begin{center}
  \includegraphics[width=0.4\linewidth]{distanceSmallTable}\hfill
  \includegraphics[width=0.5\linewidth]{neighbourJoin}
\end{center}
\end{slide}

%%%%%%%%%%%%%%%%%%%%%%% Next Slide %%%%%%%%%%%%%%%%%%%%%%%

\begin{slide}
\section{Broken Additivity}

\begin{PauseHighLight}
  \begin{itemize}
  \item If mutations just accumulated then additivity would follow\pause 
  \item However, backward mutations or two species finding the same
    mutation breaks additivity\pause
  \item Thus, to use additivity we must look at large parts of the genome
    where we are not going to be thrown out by chance events like
    backward mutation\pause
  \end{itemize}
\end{PauseHighLight}

\end{slide}

%%%%%%%%%%%%%%%%%%%%%%% Next Slide %%%%%%%%%%%%%%%%%%%%%%%
\Outline % Parsimony
%%%%%%%%%%%%%%%%%%%%%%% Next Slide %%%%%%%%%%%%%%%%%%%%%%%

\begin{slide}
\section{Character-Based Evolutionary Reconstruction}

\begin{PauseHighLight}
  \begin{itemize}
  \item A second major approach to tree construction is to find a tree
    that explains the evolutionary history of a set of sequences with as
    few mutations as possible\pause
  \item These are known as \emph{maximum parsimony trees}\pause
  \item They are the most commonly used trees\pause
  \item Here we are not given distances, but sequences\pause
  \end{itemize}
\end{PauseHighLight}

\end{slide}

%%%%%%%%%%%%%%%%%%%%%%% Next Slide %%%%%%%%%%%%%%%%%%%%%%%

\begin{slide}
\section{Aligned Protein Sequences}

\begin{center}
  \includegraphics[width=\linewidth]{alignedProteinSequences}
\end{center}
\end{slide}


%%%%%%%%%%%%%%%%%%%%%%% Next Slide %%%%%%%%%%%%%%%%%%%%%%%

\begin{slide}
\section[-1]{Algorithms for Parsimony}

\begin{PauseHighLight}
  \begin{itemize}
  \item There is no efficient algorithm for finding maximum
    parsimony trees\pause
  \item We have to check all possible trees (although we can
    use branch and bound)\pause
  \item This is feasible for a moderate number of sequences
    ($<50$)\pause
  \item Need to consider
    \begin{itemize}
    \item Assigning a cost to each tree\pause
    \item Iterating through all possible trees\pause
    \end{itemize}
  \end{itemize}
\end{PauseHighLight}

\end{slide}

%%%%%%%%%%%%%%%%%%%%%%% Next Slide %%%%%%%%%%%%%%%%%%%%%%%

\begin{slide}
\section{Costing Trees}

\begin{PauseHighLight}
  \begin{itemize}
  \item We assume that we have a set of aligned sequences\pause
  \item We treat each site independently!\pause
  \item The total cost is the sum of the costs at each site\pause
  \item We assume that we are given a binary tree whose leaves are the
    sequences\pause
  \item Each node will either be a leaf or will have left and right
    children\pause
  \end{itemize}
\end{PauseHighLight}

\end{slide}

%%%%%%%%%%%%%%%%%%%%%%% Next Slide %%%%%%%%%%%%%%%%%%%%%%%

\begin{slide}
\section{Costs of Sequence}

\begin{PauseHighLight}
  \begin{itemize}
  \item We assume that the element $a$ at each site belongs to an
    alphabet $\mathcal{A}$\pause
  \item For DNA this would be $\mathcal{A} =\{A, C, G, T, -\}$ where $-$
    represents gaps (insertion or deletion)\pause
  \item For proteins it would the amino acid and a gap\pause
  \item We consider ``weighted parsimony'' where we weight the cost of a
    substitution by a matrix $s(a,b)$\pause{} welcome back
    Blosum50\pauseb
  \item Let $C_T(a)$ be the cost of (sub-)tree $T$ given the sequence is
    $a$\pause 
  \end{itemize}
\end{PauseHighLight}

\end{slide}

%%%%%%%%%%%%%%%%%%%%%%% Next Slide %%%%%%%%%%%%%%%%%%%%%%%

\begin{slide}
\section{Dynamic Programming Solution}

\begin{pseudo}
  tree::cost() {
    cost = [];                                     % initialise array$\pause$
    if tree.isLeaf() {                           % leaf        
      for a$\in \mathcal{A}$ {
        cost[a] = $\infty$
      }
      cost[x] = 0;                                 % actual sequence$\pause$
    } else {
      for a$\in \mathcal{A}$ {
        leftCost = left.cost()                     % recursively find cost
        rightCost = right.cost()$\pause$
        cost[a] = $\min_b$ leftCost[b] + S(a,b)    % dynamic program
        cost[a] += $\min_b$ rightCost[b] + S(a,b)$\pause$
      }
    }
    return cost                                     % return array
  }$\pause$
\end{pseudo}


\end{slide}

%%%%%%%%%%%%%%%%%%%%%%% Next Slide %%%%%%%%%%%%%%%%%%%%%%%

\begin{slide}
\section{Computing Costs}

\hspace{2cm}\makebox(0,0){\includegraphics[width=0.2\linewidth]{parsimonySub}}
\pause
\pb
\pauselevel{=1}
\begin{center}
  \multipdf[width=\linewidth]{parsimonyCost}\pause
\end{center}

\end{slide}

%%%%%%%%%%%%%%%%%%%%%%% Next Slide %%%%%%%%%%%%%%%%%%%%%%%

\begin{slide}
\section{Costs for Sequences}

\begin{PauseHighLight}
  \begin{itemize}
  \item For sequences we compute the costs at each site\pause
    \begin{center}
      \includegraphics[width=\linewidth]{parsimonySeq}
    \end{center}
  \end{itemize}
\end{PauseHighLight}

\end{slide}

%%%%%%%%%%%%%%%%%%%%%%% Next Slide %%%%%%%%%%%%%%%%%%%%%%%

\begin{slide}
\section{Many Different Trees}

\begin{PauseHighLight}
  \begin{itemize}
  \item For a set of sequences there are many different possible trees
    with different costs\pause
    \begin{center}
      \includegraphics[width=\linewidth]{parsimonyFull}
    \end{center}
  \item The parsimony problem is to find the tree with the smallest cost\pause
  \end{itemize}
\end{PauseHighLight}

\end{slide}

%%%%%%%%%%%%%%%%%%%%%%% Next Slide %%%%%%%%%%%%%%%%%%%%%%%

\begin{slide}
\section[-2]{Enumerating Trees}

\pause
\pb
\pauselevel{=1}
\begin{center}
  \multipdf[width=0.9\linewidth]{binTreeGen}\pause
\end{center}

\end{slide}

%%%%%%%%%%%%%%%%%%%%%%% Next Slide %%%%%%%%%%%%%%%%%%%%%%%

\begin{slide}
\section{Branch and Bound}

\begin{PauseHighLight}
  \begin{itemize}
  \item Note that in the above procedure we generated partial
    trees\pause
  \item Any tree generated from a partial tree will have at least as
    high a cost as the partial tree\pause
  \item If the partial tree cost equals or exceeds the cost of the best
    found tree then we don't have to continue building the tree\pause
  \item This can reduce the amount of search considerable\pause
  \end{itemize}
\end{PauseHighLight}

\end{slide}
%%%%%%%%%%%%%%%%%%%%%%% Next Slide %%%%%%%%%%%%%%%%%%%%%%%

\begin{slide}
\section[-2]{Branch and Bound}

\pause
\pb
\pauselevel{=1}
\begin{center}
  \multipdf[width=0.9\linewidth]{binTreeBB}\pause
\end{center}

\end{slide}


%%%%%%%%%%%%%%%%%%%%%%% Next Slide %%%%%%%%%%%%%%%%%%%%%%%

\begin{slide}
\section[-1]{Assessing Trees}


\pause
\pb\pauselevel{=1}
\begin{itemize}
\item How reliable are the maximum parsimony trees we find?\pauseh
\item Difficult to assess directly\pauseh
\item Can build different trees by \emph{bootstrapping}\pauseh
  \begin{center}
    \multipdf[width=0.9\linewidth]{seqBootStrapping}\pause
  \end{center}
\item See if trees look the same\pause
\end{itemize}


\end{slide}

%%%%%%%%%%%%%%%%%%%%%%% Next Slide %%%%%%%%%%%%%%%%%%%%%%%

\begin{slide}
\section{Aligning and Tree Building}

\begin{PauseHighLight}
  \begin{itemize}
  \item To score the trees we want to know the alignment\pause
  \item But to do an alignment we want to find related sequences\pause
  \item It therefore makes sense to do the tree building and alignment
    together\pause
  \item There is an algorithm due to Hein that does this, but this is
    getting too involved\pause
  \end{itemize}
\end{PauseHighLight}

\end{slide}


%%%%%%%%%%%%%%%%%%%%%%% Next Slide %%%%%%%%%%%%%%%%%%%%%%%

\begin{slide}
\section{Conclusion}

\begin{PauseHighLight}
  \begin{itemize}
  \item When we have a set of sequences with a good measure of
    evolutionary distance we can use either an ultrametric model (UPGMA)
    or and additive model\pause
  \item When given a set of (aligned) sequence and some substitution
    costs then we can use maximum parsimony to find good trees\pause
  \item Uses both dynamics programming (to do the scoring) and branch
    and bound (to search through all trees)\pause
  \item These get complicated, but are scientifically important\pause
  \end{itemize}
\end{PauseHighLight}

\end{slide}
